\subsection{Limitations of evaluation design}\label{subsec:eval-limitations}

The evaluation methodology has several limitations. These limitations are listed here to bound the interpretations of the numerical outputs of the pipeline and the experiments, reported in Chapter~\ref{chapter:Results}; their implications are discussed in Chapter~\ref{chapter:Discussion}. 

\begin{itemize}
    \item \textbf{Silver labels are LLM-generated, not ground truth.} The precision of the silver labels is bounded by the biases of the LLM that generated them.
    \item \textbf{The document-extraction rubric set is small.} The document-extraction rubric set consists of 27 papers, manually and randomly sampled from the ingested corpus, and frozen for reproducibility. Since the paper set is small, the coverage of different figure-table extraction failure modes might be limited.
    \item \textbf{The synthetic relation-classification dataset is easy-skewed.} The relation-classification dataset is synthetically generated, due to the sparsity of the label classes emitted by the knowledge-extraction pipeline. The generated set of claim-pairs might consist of too cleanly related or obviously contradictory, or totally unrelated examples, which can be correctly labeled easily by the NLI model. The reported accuracy for relation classification, therefore, is an optimistic ceiling, not a corpus-representative estimate.
    \item \textbf{Cost-quality results depend on the selected voter models.} The reported escalation percentages and incurred costs highly depend on the chosen voter models and their settings.
    \item \textbf{The silver generator and the cascade's strongest voter share a model family.} The cascade's Level 3 voter is Claude Sonnet 4.6, and the silver labels are produced by Claude Opus 4.7; they are both Anthropic models. A single-Sonnet system, or the cascade with a high escalation rate, might be matching the silver not because of independent correctness, but due to family resemblance, inflating the strict-F1 scores of Sonnet and of any cascade that escalates to it. This is a validity threat for the cost-quality comparison of \textbf{RQ3}, and can only be removed by comparing the outputs to human labels.
    \item \textbf{The corpus is histopathology-focused.} The selected and ingested papers are all histopathology papers in English, selected from the PMC Open Access Subset. Therefore, the results might not be generalizable to other biomedical subfields, non-PMC papers, or non-English literature.
    \item \textbf{A single person annotated the document-extraction artifacts.} Inter-annotator agreement is not measured on the 27-paper document-extraction dataset, as the labels were produced by a single annotator under the rubric defined in Subsection~\ref{subsec:eval-doc-extraction}.
    \item \textbf{NLI four-mode ablations do not have per-pair accuracy labels.} Relation counts for each of the four modes are reported, but a comparison for their accuracy cannot be made, as the pairs do not have ground truth or silver labels.
\end{itemize}